\subsection{泛系数公式}\etiquette[1.5]{XVIII.subsec1.5}

\subsubsection{}\etiquette[1.5.1]{XVIII.subsubsec1.5.1}
设 $f:X\rightarrow S$ 是 $S$ 上的光滑射影曲线。以 $\PIC(X/S)$ 表示 $S$ 的大 fppf 景（\hyperrefext[XVII][XVII.subsec0.10]{0.10}）上的 Picard 叠，即 $X$ 上可逆层叠的正像（\hyperref[XVIII.subsubsec1.4.9]{1.4.9}）。有
\begin{align*}
\PIC(X/S)&= f_{\ast}\ch\GG_{m}[1]\tag{1.5.1.1}\\
\PIC(X/S)^{\flat} &=\tau_{\leq 0}
Rf_{\ast}(\GG_{m}[1]){\quoi}.\tag{1.5.1.2}
\end{align*}\etiquette[1.5.1.1]{XVIII.eq1.5.1.1}
\etiquette[1.5.1.2]{XVIII.eq1.5.1.2}

$X$ 的每个截面 $t$ 都在 $X$ 上定义一个可逆层 $\Oo(t)$。采用 \hyperref[XVIII.subsubsec1.4.24]{1.4.24} 的记号（\cf \hyperref[XVIII.subsubsec1.4.24]{1.4.24} $a\leftrightarrow b$），由于这一构造与任意基变换相容，它定义了 $S$ 上的一个 Picard 叠态射
\begin{equation*}
\ch \ZZ^{(X)}\longrightarrow \PIC(X/S){\quoi}.\tag{1.5.1.3}
\end{equation*}\etiquette[1.5.1.3]{XVIII.eq1.5.1.3}

若 $C$ 是 $S$ 上的任意 Picard 叠，则态射 \eqhyperref[XVIII.eq1.5.1.3]{1.5.1.3} 和 \eqhyperref[XVIII.eq1.4.24.1]{1.4.24.1} 定义一个加法函子
\begin{equation*}
\HOM(\PIC(X/S),C)\longrightarrow\HOM(\ch(\ZZ^{(X)}),
C)\xrightarrow{\sim} f_{\ast}f^{\ast}C{\quoi}.\tag{1.5.1.4}
\end{equation*}\etiquette[1.5.1.4]{XVIII.eq1.5.1.4}

应用构造 $\flat$ 后，函子 \eqhyperref[XVIII.eq1.5.1.3]{1.5.1.3} 定义一个态射\nde{这里应当理解，平移函子 $[1]$ 的运算优先级高于截断 $\tau_{\leq 0}$。}
\begin{equation*}
\ZZ^{(X)}\longrightarrow \tau_{\leq 0}
Rf_{\ast}\GG_{m}[1]\longrightarrow
Rf_{\ast}\GG_{m}[1]{\quoi},\tag{1.5.1.5}
\end{equation*}\etiquette[1.5.1.5]{XVIII.eq1.5.1.5}%
而与函子 \eqhyperref[XVIII.eq1.5.1.4]{1.5.1.4} 类似的对象，由 \eqhyperref[XVIII.eq1.4.24.2]{1.4.24.2}，是一个态射
\begin{equation*}
R\Hhom(\tau_{\leq 0} Rf_{\ast}(\GG_{m}[1]),K)\longrightarrow
Rf_{\ast} f^{\ast}K{\quoi}.\tag{1.5.1.6}
\end{equation*}\etiquette[1.5.1.6]{XVIII.eq1.5.1.6}

若 $C$ 是满足（\hyperref[XVIII.par1.3.1.1]{1.3.1.1}）的一个群 $G$ 下的挠子叠 $\ch(G[1])$，定理 \hyperref[XVIII.thm1.3.10]{1.3.10} 断言，函子 \eqhyperref[XVIII.eq1.5.1.4]{1.5.1.4} 是等价，\pageoriginale{}其逆等价将 $G_{X}$ 下的挠子 $K$ 对应到加法函子 $\langle \ast,K]$。\eqhyperref[XVIII.eq1.5.1.4]{1.5.1.4} 为等价这一断言，等价于说从 \eqhyperref[XVIII.eq1.5.1.4]{1.5.1.4} 或 \eqhyperref[XVIII.eq1.5.1.6]{1.5.1.6} 得到的态射
\begin{equation*}
\tau_{\leq 1}R\Hhom(\tau_{\leq
0}Rf_{\ast}(\GG_{m}[1]),G)\longrightarrow \tau_{\leq 1}
Rf_{\ast}f^{\ast}G\tag{1.5.1.7}
\end{equation*}\etiquette[1.5.1.7]{XVIII.eq1.5.1.7}%
是同构。

\begin{enonce*}{定理 1.5.2}\etiquette[1.5.2]{XVIII.thm1.5.2}
\emph{（泛系数公式，universal coefficient formula）。}

设 $f:X\rightarrow S$ 是 $S$ 上的光滑射影曲线，$C$ 是 $S_{\text{\rm fppf}}$ 上的 Picard 叠，并设 $K=C^{\flat}$（\hyperref[XVIII.prop1.4.15]{1.4.15}）。假设 $K$ 局部同构（在 $D^{+}(S_{\text{\rm fppf}})$ 中）于形如 $G_{-1}\rightarrow G_{0}$ 的复形，其中层 $G_{i}$ 满足（\hyperref[XVIII.par1.3.1.1]{1.3.1.1}）。则：
\begin{enumerate}
\renewcommand{\theenumi}{\Roman{enumi}}
\renewcommand{\labelenumi}{(\theenumi)}
\item 函子 \eqhyperref[XVIII.eq1.5.1.4]{1.5.1.4}
\begin{equation*}
\HOM(\PIC(X/S),C)\longrightarrow f_{\ast}f^{\ast}C\tag{1.5.2.1}
\end{equation*}\etiquette[1.5.2.1]{XVIII.eq1.5.2.1}%
是 Picard 叠的等价。

\item 从 \eqhyperref[XVIII.eq1.5.1.6]{1.5.1.6} 或 \eqhyperref[XVIII.eq1.5.1.4]{1.5.1.4} 得到的态射
\begin{equation*}
\tau_{\leq 0}R\Hhom(\tau_{\leq 0}
Rf_{\ast}(\GG_{m}[1]),K)\longrightarrow \tau_{\leq 0}
Rf_{\ast}f^{\ast}K\tag{1.5.2.2}
\end{equation*}\etiquette[1.5.2.2]{XVIII.eq1.5.2.2}%
是同构。
\end{enumerate}
\end{enonce*}

这一定理使 $\tau_{\leq 0}Rf_{\ast}(\GG_{m}[1])$ 扮演经典泛系数公式中同调所扮演的角色。

显然，对每个叠 $C$ 都有 $(I)\leftrightarrow (II)$，而且 $(I)$、$(II)$ 关于 $S$ 都是局部性质。若 $C$ 形如 $\ch(G[1])$，则如上所述，定理由 \hyperref[XVIII.thm1.3.10]{1.3.10} 得到。还需证明：若 $K$ 形如
$$
K:d:G_{-1}\longrightarrow G_{0}{\quoi},
$$
且 \eqhyperref[XVIII.eq1.5.2.2]{1.5.2.2} 对 $G_{-1}[1]$ 和 $G_{0}[1]$ 为同构，则 \eqhyperref[XVIII.eq1.5.2.2]{1.5.2.2} 对 $K$ 也是同构。为验证这一点，只需将上同调长正合列
$$
0\longrightarrow R^{-1}f_{\ast}K\longrightarrow
f_{\ast}G_{-1}\longrightarrow f_{\ast}G_{0}\longrightarrow
R^{0}f_{\ast}K\longrightarrow R^{1} f_{\ast}G_{-1}\longrightarrow
R^{1}f_{\ast}G_{0}
$$
与\pageoriginale \eqhyperref[XVIII.eq1.5.2.2]{1.5.2.2} 左侧的类似长正合列比较，并应用五引理。

\setcounter{subsubsection}{2}
\subsubsection{}\etiquette[1.5.3]{XVIII.subsubsec1.5.3}
设 $u$ 是 $S$ 上光滑射影曲线之间的一个平坦态射
$$
\xymatrix@=1.25cm{
X              \ar[rr]^{u} \ar[rd]_{f_1}&       &Y \ar[ld]^{f_2}\\
                                       &  S    &           .
}
$$
于是，加法函子图
$$\vcenter{\etiquette[1.5.3.1]{XVIII.eq1.5.3.1}
\xymatrix{
\ZZ^{(X)}  \ar[r]\ar[d]^{u} & \PIC(X/S) \ar[d]^{N_{X/Y}}\\
\ZZ^{(Y)}  \ar[r] & \PIC(Y/S)
}}\leqno{(1.5.3.1)}
$$
本质交换。

因此，对于 $S$ 上的每个 Picard 叠 $C$，图
$$\vcenter{\etiquette[1.5.3.2]{XVIII.eq1.5.3.2}
\xymatrix@R=1.75cm@C=2cm{
\HOM(\PIC(Y/S),C)         \ar[r] \ar[d]^{\HOM(N_{X/Y},C)}&
f_{2\ast}f^{\ast}_2 C \ar[d]^{u^{\ast}}\\
\HOM(\PIC(X/S), C)        \ar[r] & f_{1\ast}f^{\ast}_1 C
}}\leqno{(1.5.3.2)}
$$
本质交换；对于 $K\in \Ob D^{+}(S)$，图
$$\vcenter{\etiquette[1.5.3.3]{XVIII.eq1.5.3.3}
\xymatrix@C=2cm@R=1.5cm{
R\Hhom(\tau_{\leq 0} Rf_{2\ast}\GG_m[1],K) \ar[r]\ar[d]^{R\Hhom(N_{X/Y},K)}&
Rf_{2\ast}f^{\ast}_2 K  \ar[d]^{u^{\ast}}\\
R\Hhom(\tau_{\leq 0} Rf_{1\ast}\GG_m[1],K) \ar[r]&
Rf_{1\ast}f^{\ast}_{1} K
}}\leqno{(1.5.3.3)}
$$
交换。

再设 $G$ 是 $S$ 上满足（\hyperref[XVIII.par1.3.1.1]{1.3.1.1}）的阿贝尔层。迹 \hyperrefext[XVII][XVII.subsec6.3]{6.3} 给出一个加法函子
\begin{equation*}
\Tr_{u}:f_{1\ast}f^{\ast}_{1}\ch(G[1])\longrightarrow f_{2\ast}f^{\ast}_{2}\ch(G[1])\tag{1.5.3.4}
\end{equation*}\etiquette[1.5.3.4]{XVIII.eq1.5.3.4}%
而\pageoriginale \hyperref[XVIII.par1.3.8.6]{1.3.8.6} 给出一个使下图本质交换的同构
$$\vcenter{\etiquette[1.5.3.5]{XVIII.eq1.5.3.5}
\xymatrix@R=1.75cm{
\HOM(\PIC(X/S), \ch(G[1])) \ar[d]^{\HOM(u^{\ast}, \ch(G[1]))}&  f_{1\ast}f^{\ast}_1
\ch(G[1]) \ar[l] \ar[d]^{\Tr_u}\\
\HOM(\PIC(Y/S), \ch(G[1])) & f_{2\ast}f^{\ast}_2 \ch(G[1]).\ar[l]
}}\leqno{(1.5.3.5)}
$$

因此，图
$$\vcenter{\etiquette[1.5.3.6]{XVIII.eq1.5.3.6}
\xymatrix@R=1.75cm{
\tau_{\leq 1} R\Hhom(\tau_{\leq 0} Rf_{1\ast} \GG_{m}[1],
G)  \ar[d]_{R\Hhom(u^{\ast},G)}  & \tau_{\leq 1} Rf_{1\ast}
f^{\ast}_1G  \ar[l]_-{\sim} \ar[d]^{Tr_u}\\
\tau_{\leq 1} R\Hhom(\tau_{\leq 0} Rf_{2\ast} \GG_m[1], G)&
  \tau_{\leq 1} Rf_{2\ast}f^{\ast}_2 G \ar[l]_-{\sim}
}}\leqno{(1.5.3.6)}
$$
交换。

\subsubsection{}\etiquette[1.5.4]{XVIII.subsubsec1.5.4}
采用 \hyperref[XVIII.thm1.5.2]{1.5.2} 的记号，设 $G$ 是 $S$ 上满足 \hyperref[XVIII.par1.3.1.1]{1.3.1.1} 的阿贝尔层。关于 $S$ 对平展拓扑局部地，有（非典范地）
$$
\tau_{\leq 0}
Rf_{\ast}\GG_{m}[1]\simeq
f_{\ast}\GG_{m}[1]+R^{1}f_{\ast}\GG_{m}{\quoi}.
$$

事实上，在 Picard 叠的层次上，可将可逆层的每个同构类对应到属于该类、且沿 $X/S$ 的适当截面刚化的可逆层，由此定义一个截面 $\ch(R^{1}f_{\ast}\GG_{m})\rightarrow f_{\ast}\ch(\GG_{m}[1])$\nde{回顾此叠等同于 $\PIC(X/S)$，\cf \eqhyperref[XVIII.eq1.5.1.1]{1.5.1.1}。}。

因此，同构 \eqhyperref[XVIII.eq1.5.2.2]{1.5.2.2} 给出一个同构
\begin{equation*}
\Hhom(R^{1}f_{\ast}\GG_{m},G)\xrightarrow{\sim}\Hhom_{S}(X,G){\quoi}\tag{1.5.4.1}
\end{equation*}\etiquette[1.5.4.1]{XVIII.eq1.5.4.1}%
和一个正合列
\begin{equation*}
0\longrightarrow
\Eext^{1}(R^{1}f_{\ast}\GG_{m},G)\longrightarrow
R^1f_{\ast}G\longrightarrow
\Hhom(f_{\ast}\GG_{m},G)\longrightarrow 0{\quoi}.\tag{1.5.4.2}
\end{equation*}\etiquette[1.5.4.2]{XVIII.eq1.5.4.2}

\eqhyperref[XVIII.eq1.5.4.1]{1.5.4.1} 和 \eqhyperref[XVIII.eq1.5.4.2]{1.5.4.2} 中的态射解释如下：
\begin{enumerate}
\renewcommand{\theenumi}{\alph{enumi}}
\renewcommand{\labelenumi}{\theenumi)}
\item 态射\pageoriginale \eqhyperref[XVIII.eq1.5.4.1]{1.5.4.1} 和 \eqhyperref[XVIII.eq1.5.4.2]{1.5.4.2} 中的第一个态射，由从 $X$ 到 $\sheafPic_{X/S}$ 的典范映射定义。

\item \eqhyperref[XVIII.eq1.5.4.2]{1.5.4.2} 中的第二个态射是次数 \eqhyperref[XVIII.eq1.3.7.2]{1.3.7.2}。
\end{enumerate}

关于 $S$ 对平展拓扑局部地，有（非典范地）$\sheafPic_{X/S}\sim \sheafPic^{0}_{X/S}\times \ZZ^{k}$；所以
$$
\Eext^{1}(R^{1}f_{\ast}\GG_{m},G)\xrightarrow{\sim}\Eext^{1}(\sheafPic^{0}_{X/S}, G)
$$
而 \eqhyperref[XVIII.eq1.5.4.2]{1.5.4.2} 也可以写成
\begin{equation*}
0 \longrightarrow \Eext^{1}(\sheafPic^{0}_{X/S},G)\longrightarrow
R^{1}f_{\ast}G\longrightarrow
\Hhom(f_{\ast}\GG_{m},G)\longrightarrow 0{\quoi}.\tag{1.5.4.3}
\end{equation*}\etiquette[1.5.4.3]{XVIII.eq1.5.4.3}

\subsubsection{}\etiquette[1.5.5]{XVIII.subsubsec1.5.5}
在 \eqhyperref[XVIII.eq1.5.2.2]{1.5.2.2} 中取 $K=\GG_{m}[1]$ 时，定理 \hyperref[XVIII.thm1.5.2]{1.5.2} 就成为 $\tau_{\leq 0}(Rf_{\ast}\GG_{m}[1])$ 的一个自对偶定理。代表 $\tau_{\leq 0} Rf_{\ast} \GG_{m}[1]$ 的任意复形，都有一个典范的三步滤过，其逐次商在唯一拟同构意义下为
$$
\sheafPic_{X/S}/\sheafPic^{0}_{X/S}\quav \sheafPic^{0}_{X/S}\quav
(f_{\ast}\GG_{m})[1]{\quoi}.
$$

前述自对偶对应于 $\sheafPic_{X/S}/\sheafPic^{0}_{X/S}$ 与 $f_{\ast}\GG_{m}$ 之间的一个对偶（取值于 $\GG_{m}$），以及 $\sheafPic^{0}_{X/S}$ 上的一个自对偶（取值于 $\GG_{m}[1]$）（Jacobian 的自对偶）。

\subsubsection{}\etiquette[1.5.6]{XVIII.subsubsec1.5.6}
对于不假定在 $S$ 上固有的光滑曲线 $S:f:X\rightarrow S$，情形明显更坏。加法群的表现似乎无法控制。不过，当 $G$ 是挠性与 $S$ 的剩余特征互素的平展挠层时，前面的论证仍可推广到如下光滑曲线：它们由 $S$ 上的固有曲线去掉一个在 $S$ 上有限的部分得到。可以期望，平展 $p$-挠层也有合理的表现；见 \textsc{Serre} \cite[Ch.~VI nᵒ 11 p.~126]{sga4.exp.XVIII.bib.6}。

这里，定理 \hyperref[XVIII.thm1.2.2]{1.2.2} 的作用将由无环性定理\pageoriginale \hyperrefext[XV][XV.cor2.2]{2.2} 取代；后者有如下推论。

\begin{enonce*}{引理 1.5.7}\etiquette[1.5.7]{XVIII.lem1.5.7}
设 $p:X\rightarrow S$ 是 $S$ 上标准仿射空间的投影，$G$ 是 $S_{\text{\rm et}}$ 上的阿贝尔挠层，其挠性与 $S$ 的剩余特征互素。仍以 $G$ 表示 $G$ 在景 $S_{\text{\rm fppf}}$ 上的逆像。则函子 $p^{\ast}$ 在 $S_{\text{\rm et}}$ 上（或 $S_{\text{\rm fppf}}$ 上，由 {\textsc{Grothendieck} \cite[Corollaire~11.9]{sga4.exp.XVIII.bib.2}}，二者等价）的 $G$ 下的挠子范畴与 $X$ 上的 $G_{X}$ 下的挠子范畴之间给出范畴等价。
\end{enonce*}

\setcounter{subsubsection}{7}
\subsubsection{}\etiquette[1.5.8]{XVIII.subsubsec1.5.8}
设 $\bar{f}:\bar{X}\rightarrow S$ 是有限表现的固有态射，其纤维均为纯一维；$Y$ 是 $\bar{X}$ 的一个在 $S$ 上有限的闭子概形，由理想 $\mathfrak{m}$ 定义，并假定 $X=\bar{X}-Y$ 是 $S$ 上的光滑曲线。另设 $G$ 是 $S_{\text{\rm et}}$ 上的挠层，其挠性与 $S$ 的剩余特征互素。我们仍以 $G$ 表示作为 $G$ 的逆像得到的 $S$ 的大 fppf 景（\hyperrefext[XVII][XVII.subsec0.10]{0.10}）上的层。

设 $\GG_{m}(\mathfrak{m})$ 是 $\GG_{m\bar{X}}$ 的子层，由在 $Y$ 上等于 1 的截面组成。$\GG_{m}(\mathfrak{m})$ 下的一个挠子，等同于 $\bar{X}$ 上的一个在 $Y$ 上被平凡化的可逆层。$X$ 上在 $S$ 上有限的每个相对除子，

都定义这样的一个可逆层。关于 $S$ 对平展拓扑局部地，存在这样的除子，给出相对丰沛可逆层。

\begin{enonce*}{引理 1.5.9}\etiquette[1.5.9]{XVIII.lem1.5.9}
在 \hyperref[XVIII.subsubsec1.5.8]{1.5.8} 的假设下，设 $S$ 仿射\nde{实际上，必须假定 $S$ 是闭点为 $s$ 的局部概形。若要保留原来的假设，则在结论中应要求 $m$ 只在 $s$ 在 $S$ 中的一个开邻域的逆像上有定义，如下文 \hyperref[XVIII.par1.5.10.1]{1.5.10.1} 中的情形。}；设 $\Oo_{\bar{X}}(1)$ 是 $\bar{X}$ 上沿 $Y$ 被平凡化的相对丰沛可逆层，$\Mm$ 是 $\bar{X}$ 上沿 $Y$ 被平凡化的可逆层。对于每一点 $s\in S$，存在整数 $n_{0}$，使得对每个整数 $n\geq n_{0}$ 和纤维 $\bar{X}_{s}$ 上 $\Mm(n)\otimes k(s)$ 的每个在 $Y_{s}$ 上等于 $1$ 的截面 $m_{s}$，都存在 $\Mm(n)$ 的一个在 $Y$ 上等于 $1$ 的截面 $m$，提升 $m_{s}$。
\end{enonce*}

归结为 $S$ 是 Noether 概形的情形（因为一族满态射在基变换之后仍是一族满态射\nde{去掉这一 Noether 假设，比乍看起来要微妙一些，因为问题是仿射的，而非线性的。若 $\bar{X}$、$Y$、$\Mm$ 是由 Noether 概形 $S_0$ 上的相应数据 $\bar{X}_0$、$Y_0$、$\Mm_0$ 通过基变换 $S\to S_0$ 得到的（假定这是局部概形之间的局部态射）（\cf \EGA IV 8），且已知所需断言在这一 Noether 情形成立，则可通过一次令人不快的图追踪得到 $S$ 上所需的结果；在追踪过程中，可能需要考虑所需断言的一个变体，要求截面在 $Y_s$ 或 $Y$ 上等于 $0$，而不是 $1$。}）。设 $\Nn$ 为 $\Mm$ 中在\pageoriginale $\bar{X}_{s}$ 和 $Y$ 上消失的局部截面组成的凝聚层。对于每个 $n$，问题的局部解组成 $\Nn(n)$ 下的一个主齐性空间，因此提升 $m_{s}$ 的障碍属于 $H^{1}(\bar{X},\Nn(n))=H^{0}(S,R^{1}\bar{f}_{\ast}\Nn(n))$。但当 $n$ 充分大时，$R^{1}\bar{f}_{\ast}\Nn(n)=0$（\EGA III 2.2.1）。

\setcounter{subsubsection}{9}
\subsubsection{}\etiquette[1.5.10]{XVIII.subsubsec1.5.10}
设 $\Ll$ 是 $\bar{X}$ 上沿 $Y$ 被平凡化的可逆层，并设 $s\in S$。考虑以下条件

\paragraph{}\etiquette[1.5.10.1]{XVIII.par1.5.10.1}
$\bar{X}_{s}$ 上 $\Ll\otimes k(s)$ 的每个在 $Y_{s}$ 上等于 1 的截面，都可提升为 $s$ 的一个邻域的逆像 $U$ 上 $\Ll$ 的一个截面，它在 $\fprod{Y}{U}{S}$ 上等于 1。

\paragraph{}\etiquette[1.5.10.2]{XVIII.par1.5.10.2}
存在 $\bar{X}_{s}$ 上 $\Ll\otimes k(s)$ 的四个截面 $m_{i}$，它们在 $Y_s$ 上等于 1，且在 $X_{s}$ 的每个不可约分支的一般点线性无关。

考虑下面的条件，它们分别涉及 $\Ll$ 的一个截面 $m$、一对截面 $(m_{1},m_{2})$ 及 $\Ll$ 的三个截面 $(m_{1},m_{2},m_{3})$。

\paragraph{}\etiquette[1.5.10.3]{XVIII.par1.5.10.3}
$m$ 在 $Y$ 上等于 1，且由 $m=0$ 定义的 $X$ 的子概形 $Z(m)$ 是 $S$ 上的一个相对除子（这意味着它在 $S$ 上拟有限，因而在 $S$ 上有限，因为 $X=\bar{X}-Y$ 在 $S$ 上光滑）。

\paragraph{}\etiquette[1.5.10.4]{XVIII.par1.5.10.4}
设 $\lambda$ 为仿射直线 $g:\EE^{1}_{S}\rightarrow S$ 上的典范坐标。${g'}^{\ast}\Ll$ 的截面 $\lambda m_{1}+(1-\lambda)m_{2}$ 满足（\hyperref[XVIII.par1.5.10.3]{1.5.10.3}）。

\paragraph{}\etiquette[1.5.10.5]{XVIII.par1.5.10.5}
类似地，$\lambda m_{1}+\mu m_{2}+(1-\lambda -\mu)m_{3}$ 满足（\hyperref[XVIII.par1.5.10.3]{1.5.10.3}）。

下面的引理留给读者用一般位置论证来验证\nde{这里应注意，考虑到（\hyperref[XVIII.par1.5.10.1]{1.5.10.1}）以及在局部情形只考察 $s$ 上的纤维就能验证条件（\hyperref[XVIII.par1.5.10.3]{1.5.10.3}）的第二部分这一事实，引理现在只涉及 $\Ll$ 在纤维 $X_s$ 上的限制的截面。例如，对于 (I)，条件（\hyperref[XVIII.par1.5.10.2]{1.5.10.2}）给出一个维数为 $3$ 的\emph{仿射}空间 $V$，其中的元素是 $X_s$ 上 $\Ll\otimes k(s)$ 的在 $Y_s$ 上等于 $1$ 的截面，而且属于这个空间 $V$ 的截面在 $X_s$ 的任何一般点都不为零。我们要确定 $k\in V$，使对每个 $i$，$(m_i,k)$ 满足（\hyperref[XVIII.par1.5.10.4]{1.5.10.4}），即使连接点 $m_i$ 与 $k$ 的仿射直线不经过原点（在原点与 $X_s$ 的不同不可约分支相联系的各个表现中均如此）；换言之，只需 $k$ 不属于连接 $m_i$ 与原点的直线。于是得到一组同时要求的条件，它们分别在 $V$ 的非空开子集上满足。为了完全严格，并考虑到不可约分支未必几何不可约这一事实，当必须在一个扩张中取 $V$ 的点时（若 $k(s)$ 为有限域），实际上需要引入一个有限 Galois 扩张 $s'\to s$，使得 $X_{s'}$ 的不可约分支几何不可约。在 $X_{s'}$ 上作上述论证得到的 $V_{s'}$ 的稠密开子集，是 Galois 不变的，因而定义了所需的 $V$ 的稠密开子集。}。

\begin{enonce*}{引理 1.5.11}\etiquette[1.5.11]{XVIII.lem1.5.11}
设 $\Ll$ 是 $\bar{X}$ 上沿 $Y$ 被平凡化的可逆层，$\bar{s}$ 是 $S$ 的一个几何点，其在 $S$ 中的像为 $s$。若 $\Ll$ 在 $s$ 满足（\hyperref[XVIII.par1.5.10.1]{1.5.10.1}）和（\hyperref[XVIII.par1.5.10.2]{1.5.10.2}），则：
\begin{enumerate}
\renewcommand{\theenumi}{\Roman{enumi}}
\renewcommand{\labelenumi}{(\theenumi)}
\item 若 $\Ll$ 的 $n$ 个截面 $m_{i}$ 满足（\hyperref[XVIII.par1.5.10.3]{1.5.10.3}），则存在 $\bar{s}$ 的一个平展邻域 $U$ 和 $\Ll|\bar{f}^{-1}(U)$ 的一个截面 $k$，使各对 $(m_{i},k)$ 满足（\hyperref[XVIII.par1.5.10.4]{1.5.10.4}）。

\item 设\pageoriginale $m_{1}$、$m_{2}$、$k_{1}$、$k_{2}$ 是 $\Ll$ 的四个截面，使各对 $(m_{i},k_{j})$ 满足（\hyperref[XVIII.par1.5.10.4]{1.5.10.4}）。则存在 $\bar{s}$ 的一个平展邻域 $U$ 和 $\Ll|\bar{f}^{-1}(U)$ 的一个截面 $k$，使各三元组 $(m_{i}, k_{j}, k)$ 满足（\hyperref[XVIII.par1.5.10.5]{1.5.10.5}）。
\end{enumerate}
\end{enonce*}

当 $k(s)$ 是无限域时，只需取 $U$ 为一个 Zariski 邻域即可。

\setcounter{subsubsection}{11}
\subsubsection{}\etiquette[1.5.12]{XVIII.subsubsec1.5.12}
再设 $\Ll$ 是 $\bar{X}$ 上沿 $Y$ 刚化的可逆层，$K$ 是 $X$ 上的一个 $G$-挠子，其中 $G$ 如 \hyperref[XVIII.lem1.5.7]{1.5.7} 中所述。我们要在 $S$ 上定义一个 $G$-挠子 $\langle \Ll,K]$，将 \hyperref[XVIII.subsec1.3]{1.3} 中所研究的挠子推广，并使之关于 $\Ll$ 具有加法性。

设 $(m_{1},m_{2})$ 是 $\Ll$ 的一对满足（\hyperref[XVIII.par1.5.10.4]{1.5.10.4}）的截面。采用 \hyperref[XVIII.par1.5.10.3]{1.5.10.3} 及 \hyperref[XVIII.par1.5.10.4]{1.5.10.4} 的记号，将挠子 ${g'}^{\ast}K$ 从 $\lambda m_{1}+(1-\lambda)m_{2}=0$ 取迹到 $\EE^{1}_{S}$，得到 $\EE^{1}_{S}$ 上的一个挠子 $K_{1,2}$。它沿截面 $\lambda =0$（\resp $\lambda = 1$）的拉回，是挠子 $K_{2}=\Tr_{Z(m_{2})/S}(K)$（\resp $K_{1}=\Tr_{Z(m_{1})/S}(K)$）。由 \hyperref[XVIII.lem1.5.7]{1.5.7}，挠子 $K_{1,2}$ 典范地是 $S$ 上某个 $G$ 下的挠子的拉回；因此 $K_{1}$ 与 $K_{2}$ 之间有一个典范同构。

若 $m_{1}$、$m_{2}$、$k$ 是 $\Ll$ 的三个截面，使各对 $(m_{i},k)$ 满足（\hyperref[XVIII.par1.5.10.4]{1.5.10.4}），则通过复合仍可得到同构
\begin{equation*}
\Psi_{k}:K_{1}=\Tr_{Z(m_{1})/S}(K)\xrightarrow{\sim}
K_{2}=\Tr_{Z(m_{2})/S}(K){\quoi}.\tag{1.5.12.1}
\end{equation*}\etiquette[1.5.12.1]{XVIII.eq1.5.12.1}

若 $E$ 是 $X$ 上在 $S$ 上有限的相对除子，则 $\Ll\otimes \Oo(E)$ 的截面 $m'_{i}=m_{i}\otimes 1$ 仍满足（\hyperref[XVIII.par1.5.10.3]{1.5.10.3}），且有（对 $i=1,2$）
\begin{equation*}
K'_{i}=\Tr_{Z(m'_{i})/S}(K)\simeq K_{i}+\Tr_{E/S}(K)\tag{1.5.12.2}
\end{equation*}\etiquette[1.5.12.2]{XVIII.eq1.5.12.2}%
而同构 $\Psi_{k\otimes 1}:K'_{1}\rightarrow K'_{2}$ 从 $\Psi_{k}$ 得到：
\begin{equation*}
\Psi_{k\otimes 1}=\Psi_{k}+\id_{\Tr_{E/S}}(K){\quoi}.\tag{1.5.12.3}
\end{equation*}\etiquette[1.5.12.3]{XVIII.eq1.5.12.3}

为验证 $\Psi_{k}$ 不依赖于 $k$，只需在 $S$ 的每一点 $s$ 的一个平展邻域中验证；而由 \hyperref[XVIII.eq1.5.12.3]{1.5.12.3} 及 \hyperref[XVIII.lem1.5.9]{1.5.9}，我们\pageoriginale{}归结为条件（\hyperref[XVIII.par1.5.10.1]{1.5.10.1}）和（\hyperref[XVIII.par1.5.10.2]{1.5.10.2}）成立的情形。采用 \hyperref[XVIII.lem1.5.11]{1.5.11} (II) 的记号，$\Psi_{\lambda k_{i}+(1-\lambda)k}-\Psi_{k_{i}}$ 是从 $\EE^{1}_{S}$ 到 $G$ 的一个态射，在 $\lambda =1$ 处为零，因而恒为零；所以 $\Psi_{k_{i}}=\Psi_{k}$，进而 $\Psi_{k_{1}}=\Psi_{k_{2}}$。

设 $m_{1}$ 和 $m_{2}$ 是 $\Ll$ 的两个满足 \hyperref[XVIII.par1.5.10.3]{1.5.10.3} 的截面，且设 $K_{i}=\Tr_{Z(m_{i})/S}(K)$。为定义同构 $\Psi_{m_{2},m_{1}}:K_{1}\rightarrow K_{2}$，只需在 $S$ 的每一点 $s$ 的一个平展邻域内，以与基变换相容的方式定义它。若 $E$ 是 $X$ 上在 $S$ 上有限的相对除子，使得 $\Ll\otimes \Oo(E)$ 满足（\hyperref[XVIII.par1.5.10.1]{1.5.10.1}）和（\hyperref[XVIII.par1.5.10.2]{1.5.10.2}）（见 \hyperref[XVIII.lem1.5.9]{1.5.9}），且各对 $(m_{i}\otimes 1,k)$ 满足（\hyperref[XVIII.par1.5.10.4]{1.5.10.4}），则置（\cf \eqhyperref[XVIII.eq1.5.12.3]{1.5.12.3}）
$$
\Psi_{m_{2},m_{1}}+\Tr_{E/S}(K)=\Psi_{k}{\quoi}.
$$
这就定义了 $\Psi_{m_{2},m_{1}}$，而利用（\hyperref[XVIII.lem1.5.11]{1.5.11}. (I)）可验证 $\Psi_{m_{3},m_{2}}\Psi_{m_{2},m_{1}}=\Psi_{m_{3},m_{1}}$。这些构造与基变换相容。若存在 $\Ll$ 的一个满足（\hyperref[XVIII.par1.5.10.3]{1.5.10.3}）的截面 $m$，则置
$$
\langle \Ll,K]=\Tr_{Z(m)/S}(K)
$$
由前面的论证，这一挠子在典范同构意义下不依赖于 $m$ 的选择。

在一般情形，注意为定义 $\langle \Ll,K]$，只需在 $S$ 的每一点 $s$ 的一个平展邻域中，以与基变换相容的方式定义它。若如上所述，$\Ll\otimes \Oo(E)$ 满足（\hyperref[XVIII.par1.5.10.1]{1.5.10.1}）及（\hyperref[XVIII.par1.5.10.2]{1.5.10.2}），则 $\Ll\otimes \Oo(E)$ 有满足（\hyperref[XVIII.par1.5.10.3]{1.5.10.3}）的截面，并置
$$
\langle \Ll,K]=\langle \Ll\otimes \Oo(E), K]-\Tr_{E/S}(K){\quoi}.
$$

\subsubsection{}\etiquette[1.5.13]{XVIII.subsubsec1.5.13}
在 \hyperref[XVIII.subsubsec1.5.8]{1.5.8} 的假设下，若 $x$、$y$ 是 $\bar{f}_{\ast}\GG_{m}(\mathfrak{m})$ 的两个截面，则当 $\lambda$ 在 $S$ 上仿射直线的某个具有非空纤维的开子集中取值时，$\lambda x+(1-\lambda)y$ 仍是一个截面。

这可以表述为

\paragraph{}\etiquette[1.5.13.1]{XVIII.par1.5.13.1}
“$\bar{f}_{\ast}\GG_{m}(\mathfrak{m})$ 是道路连通的。”

从仿射直线的一个具有非空纤维的开子集到 $G$ 的每个\pageoriginale{}态射都是常值的。因此，不存在从 $\bar{f}_{\ast}\GG_{m}(\mathfrak{m})$ 到 $G$ 的非平凡同态，而且 $\Ll$ 的自同构在 $\langle \Ll,K]$ 上平凡地作用。所以，挠子 $\langle \Ll,K]$ 只依赖于 $K$ 和由 $\Ll$ 定义的 $R^{1}\bar{f}_{\ast}\GG_{m}(\mathfrak{m})$ 的截面。通过整体化，由此可以对 $R^{1}\bar{f}_{\ast}\GG_{m}(\mathfrak{m})$ 的截面 $\lambda$ 定义符号 $\langle \lambda,K]$。挠子 $\langle \lambda,K]$ 关于 $\lambda$ 具有加法性，因此对应于 $R^{1}\bar{f}_{\ast}\GG_{m}(\mathfrak{m})$ 由 $G$ 的一个扩张 $e(K)$（见 \hyperref[XVIII.prop1.4.23]{1.4.23}，或 \SGA 7 VII 1.1.6 和 1.2）。

$e(K)$ 的形成与任意基变换及挠子的加法相容，因而给出一个加法函子
$$
e:\bar{f}_{\ast}\ch(G[1])\longrightarrow
\EXT(R^{1}\bar{f}_{\ast}\GG_{m}(\mathfrak{m}),G){\quoi}.
$$

设 $j$ 是从 $X$ 到 $R^{1}\bar{f}_{\ast}\GG_{m}(\mathfrak{m})$ 的典范映射，$j:t\rightarrow$（$\Oo(t)$ 的类）。沿 $j$ 的拉回是一个加法函子
$$
j^{\ast}:\EXT(R^{1}\bar{f}_{\ast}\GG_{m}(\mathfrak{m}),G)\longrightarrow
\bar{f}_{\ast}\ch(G[1])
$$
显然 $j^{\ast}e\sim \Id$。如 \hyperref[XVIII.thm1.3.10]{1.3.10} 和 \hyperref[XVIII.thm1.5.2]{1.5.2} 中一样，可验证对于 fppf 拓扑，有

\begin{enonce*}{命题 1.5.14}\etiquette[1.5.14]{XVIII.prop1.5.14}
在 \hyperref[XVIII.subsubsec1.5.8]{1.5.8} 的假设下，\hyperref[XVIII.subsubsec1.5.13]{1.5.13} 中的函子 $e$、$j^{\ast}$ 是互为拟逆的等价。由 \hyperref[XVIII.const1.4.19]{1.4.19} 和 \hyperref[XVIII.prop1.4.23]{1.4.23}，得到一个同构
$$
\Ext^{1}(R^{1}\bar{f}_{\ast}\GG_{m}(\mathfrak{m}),G)\xrightarrow{\sim}
H^{1}(X,G){\quoi}.
$$
\end{enonce*}

\setcounter{subsubsection}{14}
\subsubsection{}\etiquette[1.5.15]{XVIII.subsubsec1.5.15}
设 $f_{i}:X_{i}\rightarrow S$ $(i=1,2)$ 是两条如 \hyperref[XVIII.subsubsec1.5.8]{1.5.8} 中所述的已紧化曲线。

设 $\bar{u}:\bar{X}_{1}\rightarrow \bar{X}_{2}$ 是一个态射，使得 $Y_{1}$ 是 $Y_{2}$ 的逆像的子概形，且它诱导一个从 $X_{1}$ 到 $X_{2}$ 的有限平坦态射。则对 $\bar{X}_{2}$ 上沿 $Y_2$ 刚化的可逆层 $\Ll$ 及 $X_{1}$ 上的挠子 $K$，有
$$
\langle \Ll,\tr_{X_{1}/X_{2}}(K)]\xrightarrow{\sim}
\langle \bar{u}^{\ast}\Ll,K]{\quoi},
$$
由此\pageoriginale{}推出下图交换
$$\vcenter{\etiquette[1.5.15.1]{XVIII.eq1.5.15.1}
\xymatrix@R=1.75cm{
\Ext^{1}(R^1\bar{f}_{1\ast} \GG_m (\mathfrak{m}_1), G) 
\ar[r]^-{\sim}\ar[d]^{\Ext^{1}(\bar{u}^{\ast},\Id)} &   H^{1}(X_1,G) \ar[d]^{\Tr_{X_1 /X_2}}\\
\Ext^1(R^1\bar{f}_{2\ast}\GG_m (\mathfrak{m}_2), G)  \ar[r]^-{\sim}&
H^1(X_2,G).
}}\leqno{(1.5.15.1)}
$$

若 $\bar{u}:\bar{X}_{1}\rightarrow \bar{X}_{2}$ 有限平坦，且满足 $Y_{1}=\bar{u}^{-1}Y_{2}$（作为概形），则对 $\bar{X}_{1}$ 上沿 $Y_1$ 刚化的可逆层 $\Ll$ 及 $X_{2}$ 上的挠子 $K$，有
$$
\langle \Ll,u^{\ast}K]\xrightarrow{\sim}
\langle N_{\bar{X}_{1}/\bar{X}_{2}}\Ll,K]
$$
由此推出下图交换
\begin{equation*}
\vcenter{
\xymatrix@R=1.75cm@C=1.5cm{
\Ext^{1}(R^{1}\bar{f}_{2\ast} \GG_m (\mathfrak{m}_2) , G)
\ar[d]^{\Ext^{1}(N_{\bar{X}_1 /\bar{X}_2}, \Id)}\ar[r]^-{\sim} &
H^{1}(X_2,G) \ar[d]^{u^{\ast}}\\
\Ext^{1}(R^{1}\bar{f}_{1\ast} \GG_m (\mathfrak{m}_1), G)  \ar[r]^-{\sim} &
H^{1}(X_1, G).
}}\tag{1.5.15.2}
\end{equation*}\etiquette[1.5.15.2]{XVIII.eq1.5.15.2}

当然，可以用 Picard 叠的语言更精确地陈述这一结果。

\subsection{一个消没定理}\etiquette[1.6]{XVIII.subsec1.6}

本节从 \hyperref[XVIII.lem1.6.6]{1.6.6}（第二种证明）到末尾，可以独立于第 \hyperref[XVIII.subsec1.2]{1.2} 至 \hyperref[XVIII.subsec1.5]{1.5} 小节阅读。

\begin{enonce*}{引理 1.6.1}\etiquette[1.6.1]{XVIII.lem1.6.1}
设 $X$ 是概形，$Y$ 是 $X$ 的一个由理想 $\mathfrak{m}$ 定义的闭子概形，其补集为 $U$。
$$
U\mathop{\hookrightarrow}^{j} X\mathop{\hookleftarrow}^{i} Y{\quoi}.
$$
设 $\GG_{m}(\mathfrak{m})$ 是如下 \emph{fpqc} 层：它在 $X$-概形 $X_{1}$ 上的截面，是 $\Oo^{\ast}_{X_{1}}$ 中在 $Y$ 的逆像上等于 $1$ 的截面。对在 $X$ 上可逆的 $n\geq 1$，序列
$$
0\rightarrow j_{!}𝛍_{n}\rightarrow
\GG_{m}(\mathfrak{m})\xrightarrow{n} \GG_{m}(\mathfrak{m})\rightarrow
0
$$
是 $X$ 的大平展景\emph{（\hyperrefext[XVII][XVII.subsec0.10]{0.10}）}上的层的正合列。
\end{enonce*}

设\pageoriginale $i$ 为由 $i$ 诱导的大平展景态射 $i:Y\rightarrow X$：
$$
i^{\ast}(X_{1}/X)=\fprod{X_{1}}{Y}{X}{\quoi}.
$$
函子 $i_{\ast}$ 是正合的。在图
$$\vcenter{\etiquette[1.6.1.1]{XVIII.eq1.6.1.1}
\xymatrix{
             &    0  \ar[d]           &  0  \ar[d]                   &    0\ar[d]    \\
0  \ar[r]    &j_!𝛍_n  \ar[r]\ar[d]  & \GG_m(\mathfrak{m})\ar[r]\ar[d]  &
             \GG_m(\mathfrak{m}) \ar[r] \ar[d] &0\\
0  \ar[r]     &𝛍_n  \ar[r]  \ar[d]  & \GG_m  \ar[r] \ar[d]  &
  \GG_m  \ar[r] \ar[d]  & 0\\
0  \ar[r]    & i_{\ast}𝛍_n  \ar[d] \ar[r] &
  i_{\ast}\GG_m  \ar[r]\ar[d]  & i_{\ast}\GG_m   \ar[r]\ar[d] & 0\\
             &    0             &  0                    &    0   \\
}}\leqno{(1.6.1.1)}
$$
中，各列均正合，最后两行也正合（\hyperrefext[IX][IX.Thr3.2]{3.2}）。因此第一行正合。

\setcounter{subsubsection}{1}
\subsubsection{}\etiquette[1.6.2]{XVIII.subsubsec1.6.2}
设 $k$ 是一个域，
$$
G'_{1}\longrightarrow G'_{2}\longrightarrow G\longrightarrow
G''_{2}\longrightarrow G''_{1}
$$
是 $\Spec (k)$ 的大 fppf 景上的阿贝尔层的正合列。回顾，若 $G'_{1}$、$G'_{2}$ 可表且有限型，而 $G''_{2}$、$G''_{1}$ 可表且局部有限型，则 $G$ 可表且局部有限型\acorriger{当 $G'_1=G'_2=0$ 时，这是显然的。当 $G''_1=G''_2=0$ 时，利用 \SGA 3 V 10.1.1（新版），可归结为 $G'_1\to G'_2$ 为单态射的情形，而此情形见（\SGA 3 VI${}_{\text{A}}$ 3.2）。于是我们归结为 $G'_1=G''_1=0$ 的情形。这时可将层 $G$ 看作概形 $G''_2$ 上 $G'_2$ 下的挠子。若假定 $G'_2$ 仿射（如果在进行这些化简之前假定 $G'_1$、$G'_2$ 仿射，则由（\SGA 3 VI${}_{\text{B}}$ 11.17）可归结到这种情形），则 $G$ 的可表性立即由忠实平坦下降得到。这个特殊情形足以用于本段中的应用。若 $G'_2$ 拟射影，则可利用（\SGA 1 VIII 7.7）及类似于 \cite[引理~3.3.1]{sga4.exp.XVIII.bib.4bis} 中的论证得到 $G$ 的可表性。事实上，域上的每个有限型群概形 $G$ 都拟射影：可以假定 $G$ 连通；当 $G$ 光滑时，这就是 \cite[推论~1.2]{sga4.exp.XVIII.bib.0bis}，而一般情形可以归结到光滑情形（M. Raynaud），因为由（\SGA 3 VII${}_{\text{A}}$ 8.3），存在从 $G$ 到一个光滑群概形的有限态射。}。

若 $P$ 是 $k$ 上的有限型连通交换代数群，且 $n$ 是在 $k$ 中可逆的整数，则序列
\begin{equation*}
0\rightarrow {}_{n}P\rightarrow P\xrightarrow{n} P\rightarrow 0\tag{1.6.2.1}
\end{equation*}\etiquette[1.6.2.1]{XVIII.eq1.6.2.1}%
正合。若 $G$ 是 $k$ 上被 $n$ 零化的有限平展交换群，则正合列
$$
\Hom(P,G)\rightarrow \Hom({}_{n}P,G)\xrightarrow{j}
\Ext^{1}(P,G)\xrightarrow{n} \Ext^{1}(P,G)
$$
中的第一个群和最后一个箭头为零，从而有同构
\begin{equation*}
j:\Hom({}_{n}P,G)\xrightarrow{\sim} \Ext^{1}(P,G){\quoi}.\tag{1.6.2.2}
\end{equation*}\etiquette[1.6.2.2]{XVIII.eq1.6.2.2}

\subsubsection{}\etiquette[1.6.3]{XVIII.subsubsec1.6.3}
在\pageoriginale \hyperref[XVIII.lem1.6.1]{1.6.1} 的假设下，假定 $X$ 在域 $k$ 的谱上固有
% SGA-SRC-001141
$$
p:X\longrightarrow \Spec(k){\quoi}.
$$
正合列
$$
0\longrightarrow \GG_{m}(\mathfrak{m})\longrightarrow
\GG_{m}\longrightarrow i_{\ast}\GG_{m}\longrightarrow 0
$$
定义了 $\Spec (k)$ 上 fppf 层上同调的一个正合列，容易直接将它解释为
\begin{equation*}
 0\rightarrow p_{\ast}\GG_{m}(\mathfrak{m})\rightarrow
p_{\ast}\GG_{m_{X}}\rightarrow
(pi)_{\ast}\GG_{m_{Y}}\xrightarrow{j}\tag{1.6.3.1}
 R^{1}p_{\ast}\GG_{m}(\mathfrak{m})\rightarrow
\sheafPic_{X/k}\rightarrow \sheafPic_{Y/k}{\quoi}.
\end{equation*}\etiquette[1.6.3.1]{XVIII.eq1.6.3.1}

层 $p_{\ast}\GG_{m_{X}}$ 和 $(pi)_{\ast}\GG_{m_{Y}}$ 可由有限型\sisi{}{且仿射}的群概形表出。层 $\sheafPic_{X/k}$ 和 $\sheafPic_{Y/k}$ 可由局部有限型群概形表出。由 \hyperref[XVIII.subsubsec1.6.2]{1.6.2}，层 $p_{\ast}\GG_{m}(\mathfrak{m})$ 和 $R^{1}p_{\ast}\GG_{m}(\mathfrak{m})$ 可表。由 \eqhyperref[XVIII.par1.5.13.1]{1.5.13.1}，前者由一个连通（有限型）群概形表出；后者局部有限型。

\begin{enonce*}{定义 1.6.4}\etiquette[1.6.4]{XVIII.def1.6.4}
在上述假设下，以 $\sheafPic_{\mathfrak{m},X/k}$ 表示表出 $R^{1}p_{\ast}\GG_{m}(\mathfrak{m})$ 的群概形。以 $\sheafPic^{0}_{\mathfrak{m},X/k}$ 表示其单位分支，并以 $\sheafPic^{\tau}_{\mathfrak{m},X/k}$ 表示 $\sheafPic_{\mathfrak{m},X/k}/\sheafPic^{0}_{\mathfrak{m},X/k}$ 的挠子群在 $\sheafPic_{\mathfrak{m},X/k}$ 中的逆像。
\end{enonce*}

概形 $\sheafPic_{\mathfrak{m},X/k}$ 表出由如下预层生成的 fppf 层：它将 $k$ 上的每个概形 $S$ 对应到 $\fprod{S}{X}{k}$ 上沿 $\fprod{S}{Y}{k}$ 被平凡化的可逆模的同构类集合（\hyperref[XVIII.subsubsec1.5.8]{1.5.8}）。

当 $k$ 代数闭且 $n$ 在 $k$ 中可逆时，Kummer 正合列 \hyperref[XVIII.lem1.6.1]{1.6.1} 给出一个同构
\begin{equation*}
H^{1}_{c}(U,𝛍_{n})\xrightarrow{\sim}
{}_{n}\sheafPic^{\tau}_{\mathfrak{m},X/k}(k){\quoi}.\tag{1.6.4.1}
\end{equation*}\etiquette[1.6.4.1]{XVIII.eq1.6.4.1}

\setcounter{subsubsection}{4}
\subsubsection{}\etiquette[1.6.5]{XVIII.subsubsec1.6.5}
设\pageoriginale $X$ 是代数闭域 $k$ 上的光滑曲线，$\bar{X}$ 是以 $X$ 为稠密开子概形的光滑射影曲线，$Y$ 是 $\bar{X}$ 的一个闭子概形，使得 $X=\bar{X}-Y$，$\mathfrak{m}$ 是定义 $Y$ 的理想，$n$ 是在 $k$ 中可逆的整数，$G$ 是被 $n$ 零化的有限阿贝尔群。

群 $\sheafPic_{\mathfrak{m},\bar{X}/k}$ 非典范地是子群 $\sheafPic^{0}_{\mathfrak{m},\bar{X}/k}$ 与一个自由阿贝尔群的积。因此，由 \eqhyperref[XVIII.eq1.6.4.1]{1.6.4.1}、\eqhyperref[XVIII.eq1.6.2.2]{1.6.2.2} 及 \hyperref[XVIII.prop1.5.14]{1.5.14}，有一个典范的复合同构
\begin{gather*}
\Hom(H^{1}_{c}(X,𝛍_{n}),G)\xleftarrow{\sim}
\Hom({}_{n}\sheafPic^{\tau}_{\mathfrak{m},\bar{X}/k}(k),G)=\Hom({}_{n}\sheafPic^{0}_{\mathfrak{m},\bar{X}/k},G)\\
\xrightarrow{\sim}
\Ext^{1}(\sheafPic^{0}_{\mathfrak{m},\bar{X}/k},G)\xleftarrow{\sim}
\Ext^{1}(\sheafPic_{\mathfrak{m},\bar{X}/k},G)\xrightarrow{\sim} H^{1}(X,G)
\end{gather*}

若 $u:X\rightarrow Y$ 是 $k$ 上光滑曲线之间的有限态射，则由 \hyperrefext[XVII][XVII.subsec6.3]{6.3}. 和 \eqhyperref[XVIII.eq1.5.15.2]{1.5.15.2} 得到下图交换
\begin{equation*}
\vcenter{
\xymatrix@C=2cm@R=1.5cm{
\Hom(H^{1}_c(Y,𝛍_n),G)  \ar[r]^(.55){\sim} \ar[d]^{\Hom(\Tr_u,G)} &
H^{1}(Y,G)               \ar[d]^{u^{\ast}}\\
\Hom(H^1_c(X,𝛍_n),G)    \ar[r]^(.55){\sim} & H^{1}(X,G).
}}\tag{1.6.5.1}
\end{equation*}\etiquette[1.6.5.1]{XVIII.eq1.6.5.1}

\begin{enonce*}{引理 1.6.6}\etiquette[1.6.6]{XVIII.lem1.6.6}
设 $X$ 是代数闭域 $k$ 上的光滑曲线，$n$ 是在 $k$ 中可逆的整数，$u:X'\rightarrow X$ 是一个驯分歧的平展覆盖（tame étale covering）。则态射
$$
\Tr_{u}:H^{1}_{c}(X',\ZZ/n)\longrightarrow
H^{1}_{c}(X,\ZZ/n)
$$
与态射
$$
u^{\ast}:H^{1}(X,\ZZ/n)\longrightarrow
H^{1}(X',\ZZ/n){\quoi}.
$$
的转置（非典范地\nde{若已选定一个同构 $𝛍_n\simeq \mathbf{Z}/n$，则此同构是典范的。}）同构。
\emph{第一种证明。}由于 $𝛍_{n}$ 与 $\ZZ/n$ 同构，\hyperref[XVIII.lem1.6.6]{1.6.6} 是 \eqhyperref[XVIII.eq1.6.5.1]{1.6.5.1} 中取 $G=\ZZ/n$ 的特殊情形。
\end{enonce*}

\noindent
\emph{第二种证明。}
若 $k$ 是复数域 $\CC$，则由比较定理，态射 $\Tr_{u}$ 和 $u^{\ast}$ 等同于由拓扑空间的连续映射 $u:X'(\CC)\rightarrow X(\CC)$ 所定义的类似态射：
\begin{align*}
\Tr_{u} &: H^{1}_{c}(X'(\CC),\ZZ/n) \longrightarrow
H^{1}_{c}(X(\CC), \ZZ/n)\quad \text{和}\\
u^{\ast} &: H^{1}(X(\CC),\ZZ/n)\longrightarrow H^{1}(X'(\CC),\ZZ/n){\quoi}.
\end{align*}
这\pageoriginale{}两个同态在 Poincaré 对偶下互为转置。

由 Lefschetz 原理，当 $k$ 的特征为 $0$ 时，断言 \hyperref[XVIII.lem1.6.6]{1.6.6} 仍成立（事实上，所讨论的上同调群在代数闭基域扩张下不变）。

若 $k$ 的特征为 $p>0$，设 $W(k)$ 为 $k$ 上的 Witt 向量环。设 $\bar{X}$ 是 $X$ 的光滑射影完备化。由（\SGA 1 III 7.4），曲线 $\bar{X}$ 提升为 $W(k)$ 上的一条光滑射影曲线 $\bar{X}_{1}$。由于 $\bar{X}$ 光滑，每一点 $s_{i}\in S=\bar{X}-X$ 都提升为 $\bar{X}_{1}$ 在 $W(k)$ 上的一个截面 $s'_{i}$；置 $S_{1}=⋃_{i}s'_{i}(\Spec(W(k)))$，$X_{1}=\bar{X}_{1}-S_{1}$。

回顾（\SGA 1 XIII 2.4），$X$ 的覆盖 $X'$ 提升为 $X_{1}$ 的一个平展覆盖 $X'_{1}$，而且 $X'_{1}$ 由 $W(k)$ 上的一条光滑射影曲线 $\bar{X}'_{1}$ 去掉有限个两两不交截面的并 $S'_{1}$ 得到。
$$
\xymatrix@=.5cm{
                  &     X_1 \ar[ddddd]^{f_1}       &                 &
                  X\ar[ll]\ar[ddddd]\\
 & & &\\
X'_1 \ar[ruu]_{u_1}\ar[rddd]_{f'_1}&                & X' \ar[ll]\ar[ruu]_{u} \ar[rddd] &   \\
& & &\\
& & &\\
                  &    \Spec(W(k)) &                 &\Spec(k).\ar[ll]
}
$$

平展层 $R^{1}f_{1\ast}(\ZZ/n)$、$R^{1}f'_{1\ast}(\ZZ/n)$、$R^{1}f_{1!}(\ZZ/n)$ 和 $R^{1}f'_{1!}(\ZZ/n)$ 局部常值，且其形成与基变换相容，\pageoriginale{}这由联系 $\bar{X}_{1}$、$X_{1}$、$S_{1}$ 或 $\bar{X}'_{1}$、$X'_{1}$、$S'_{1}$ 的上同调的正合列得到。

因此，态射
$$
\Tr_{u_{1}}:R^{1}f'_{1!}\ZZ/n\longrightarrow
R^{1}f_{1!}\ZZ/n
$$
或
$$
u^{\ast}_{1}:R^{1}f_{1\ast}\ZZ/n\longrightarrow
R^{1}f'_{1\ast}\ZZ/n
$$
的特殊几何纤维与一般几何纤维同构，于是归结为已经处理过的 $k$ 特征为零的情形。

\begin{enonce*}{引理 1.6.7}\etiquette[1.6.7]{XVIII.lem1.6.7}
设 $X$ 是特征指数为 $p$ 的代数闭域 $k$ 上的连通光滑曲线，$n$ 是与 $p$ 互素的整数。存在一个主覆盖\emph{（即有限平展 Galois（满射）覆盖）} $u:X'\rightarrow X$，其次数整除 $n$ 的某个幂，使得迹态射
$$
\Tr_{u}:H^{1}_{c}(X', \ZZ/n)\longrightarrow
H^{1}_{c}(X,\ZZ/n)
$$
为零。
\end{enonce*}

回顾群 $H^{1}(X,\ZZ/n)$ 是有限的（\hyperrefext[IX][IX.thm4.6]{4.6} 或 \hyperrefext[XVI][XVI.thm5.2]{5.2}）。设 $u:X'\rightarrow X$ 是 $X$ 的最大连通阿贝尔主覆盖，其 Galois 群被 $n$ 零化（其 Galois 群记作 $H_{1}(X,\ZZ/n)$）。由构造，同态
$$
u^{\ast}:H^{1}(X,\ZZ/n)\longrightarrow H^{1}(X', \ZZ/n)
$$
为零。因此，由 \hyperref[XVIII.lem1.6.6]{1.6.6}，同态
$$
\Tr_{u}:H^{1}_{c}(X',\ZZ/n)\longrightarrow
H^{1}_{c}(X,\ZZ/n)
$$
为零。

回顾下面的引理（\EGA IV 15.6.5）：

\begin{enonce*}{引理 1.6.8}\etiquette[1.6.8]{XVIII.lem1.6.8}
设 $f:X\rightarrow S$ 是光滑态射，$x$ 是 $f$ 的一个截面。则存在 $X$ 的一个包含 $x$ 的 Zariski 开子集 $U$，使得 $f|U$ 的几何纤维连通。
\end{enonce*}

\begin{enonce*}{基本引理 1.6.9}\etiquette[1.6.9]{XVIII.lemfond1.6.9}
设\pageoriginale $f:X\rightarrow S$ 是可紧化的光滑曲线（\hyperref[XVIII.def1.1.2]{1.1.2}）\emph{（\hyperrefext[XVII][XVII.subsec3.2]{3.2}）}，$x$ 是 $X$ 的一个几何点，$s$ 是其在 $S$ 中的像，$n\geq 1$ 是在 $S$ 上可逆的整数。存在 $s$ 在 $S$ 中的一个平展邻域 $V$，以及 $x$ 在 $f^{-1}(V)=\fprod{V}{X}{S}$ 中的一个平展邻域 $U$
\begin{equation*}
\vcenter{
\xymatrix@R=1.5cm@C=1cm{
x \ar[r]\ar@{-}[d]& U\ar[rd]_{f'}\ar[r]^{u} & f^{-1}(V)\ar[r] \ar[d]^{f_v}& X\ar[d]^{f}\\
s \ar[rr]&   & V\ar[r]         & S,
}}\tag{1.6.9.1}
\end{equation*}\etiquette[1.6.9.1]{XVIII.eq1.6.9.1}%
使得 $R^{0}f'_{!}\sheaf{\ZZ}/n=0$，箭头\emph{ \hyperrefext[XVII][XVII.subsec6.2]{6.2}}
$$
\Tr_{u}:R^{1}f'_{!}\sheaf{\ZZ}/n\longrightarrow
R^{1}f_{v!}\sheaf{\ZZ}/n
$$
为零，且迹态射
$$
\Tr_{f'}:R^{2}f'_{!}\sheaf{\ZZ}/n(1)\longrightarrow
\sheaf{\ZZ}/n
$$
为同构。
\end{enonce*}

为使 $R^{0}f'_{!}\sheaf{\ZZ}/n=0$，只需 $f'$ 拟仿射；为使 $\Tr_{f'}$ 为同构，只需 $f'$ 的几何纤维连通（\hyperref[XVIII.lem1.1.9]{1.1.9}）；由 \hyperref[XVIII.lem1.6.8]{1.6.8}，关于 $S$ 对平展拓扑局部地，容易获得 $f'$ 的这些性质。还需对具有几何连通纤维的 $X$，证明存在图 \eqhyperref[XVIII.eq1.6.9.1]{1.6.9.1}，使所讨论的箭头 $\Tr_{u}$ 为零：之后只需缩小 $U$、$V$，就能使 $f'$ 满足另外两个所需的性质。

问题关于 $S$ 在 $s$ 的邻域中是局部的；由 \hyperrefext[XVII][XVII.thm5.2.6]{5.2.6}（基变换），立即归结为 $S$ 是 Noether 概形的情形。通过取极限，并利用 \hyperrefext[XVII][XVII.thm5.3.6]{5.3.6}（可构造性），归结为 $S$ 是严格局部 Noether 概形，且闭点为 $s$ 的像的情形。

设 $t$ 为 $S$ 的一个几何点。由 \hyperref[XVIII.lem1.6.7]{1.6.7}，存在一个次数与 $k(t)$ 的特征指数
% SGA-SRC-001143
$p$ 互素的主覆盖 $u:X^{1}_{t}\rightarrow X_{t}$，使得箭头
$$
\Tr_{u}:H^{1}_{c}(X^{1}_{t},\ZZ/n)\longrightarrow
H^{1}_{c}(X_{t},\ZZ/n)
$$
为\pageoriginale{}零。由无环性定理 \hyperrefext[XV][XV.thm2.1]{2.1}（或 \hyperrefext[XV][XV.lem2.6]{2.6}），存在 $x$ 在 $X$ 中的一个平展邻域 $v:U\rightarrow X$，使得 $X^{1}_{t}$ 在 $U_{t}$ 上分裂，因而 $v_{t}:U_{t}\rightarrow X_{t}$ 通过 $X^{1}_{t}$ 分解。所以仍有
$$
\Tr_{v_{t}}=0:H^{1}_{c}(U_{t},\ZZ/n)\longrightarrow
H^{1}_{c}(X_{t},\ZZ/n){\quoi}.
$$

因此，对于 $S$ 的每个几何点 $t$，都存在 $x$ 在 $X$ 中的一个平展邻域
$$
\xymatrix{
U \ar[r]^{v} \ar[rd]_{f'}& X \ar[d]^{f}\\
  & S
}
$$
使得
\begin{equation*}
\Tr_{v}:R^{1}f'_{!}\ZZ/n\longrightarrow
R^{1}f_{!}\ZZ/n\tag{1.6.9.2}
\end{equation*}\etiquette[1.6.9.2]{XVIII.eq1.6.9.2}%
在 $t$ 的纤维为零（\hyperrefext[XVII][XVII.thm5.2.6]{5.2.6}）。

当 $U$ 变化时，箭头 \eqhyperref[XVIII.eq1.6.9.2]{1.6.9.2} 的像组成可构造层 $R^{1}f_{!}\ZZ/n$ 的一个滤过递减子层系统，它在 $S$ 的任意给定几何点最终变为零。利用 Noether 归纳法可验证，这样的系统总是稳定的，且稳定值为零。因此存在 $U$，使 \eqhyperref[XVIII.eq1.6.9.2]{1.6.9.2} 为零，\hyperref[XVIII.lemfond1.6.9]{1.6.9} 得证。
